\Chapter{El óptimo social y las externalidades de congestión}

\section{El problema del óptimo social}

\subsection{Formulación}

\datebox{2026}

\begin{definition}[Costo total del sistema]
El \emph{costo total del sistema} dado el vector de flujos $x \in \R^m_+$ es:
\[
\TC(x) \defeq \sum_{a \in A} x_a \cdot t_a(x_a).
\]
\end{definition}

\begin{definition}[Problema del óptimo social]\label{def:social_optimum}
El \emph{problema del óptimo social} es:
\begin{equation}\label{eq:social_optimum}
(\calSO) \quad \min_{f \in \FeasSet} \; \TC(f) \defeq \sum_{a \in A} x_a(f) \cdot t_a(x_a(f)).
\end{equation}
Una solución $f^{\SO}$ del problema $(\calSO)$ se denomina una \emph{asignación socialmente óptima}, y los flujos inducidos $x^{\SO} = \Delta f^{\SO}$ se denominan \emph{flujos del óptimo social}.
\end{definition}

\begin{proposition}\label{prop:so_convex}
El problema $(\calSO)$ es convexo. Si $2t_a'(x_a) + x_a t_a''(x_a) \geq 0$ para todo $a$ y todo $x_a \geq 0$, entonces $\TC(x)$ es convexo en $x$.
\end{proposition}

\begin{proof}
La segunda derivada de $x_a t_a(x_a)$ con respecto a $x_a$ es $2t_a'(x_a) + x_a t_a''(x_a)$. Para la función BPR:
\[
2t_a'(x_a) + x_a t_a''(x_a) = \frac{\alpha \beta t_a^0}{c_a^\beta}\left[2 x_a^{\beta-1} + (\beta-1)x_a^{\beta-1}\right] = \frac{\alpha \beta (\beta+1) t_a^0}{c_a^\beta} x_a^{\beta-1} \geq 0.
\]
Por tanto $\TC$ es convexo para la BPR.
\end{proof}

\subsection{Diferencia entre el problema de Beckmann y el óptimo social}

\begin{proposition}[No equivalencia entre UE y SO]\label{prop:ue_ne_so}
En general, $f^{\UE} \neq f^{\SO}$, es decir, el equilibrio de Wardrop no minimiza el costo total del sistema.
\end{proposition}

\begin{proof}
La función objetivo del problema de Beckmann es $Z(f) = \sum_a T_a(x_a)$, con $T_a'(x_a) = t_a(x_a)$.

La función objetivo del óptimo social es $\TC(f) = \sum_a x_a t_a(x_a)$, cuya derivada con respecto a $x_a$ es el costo marginal social $\CMS_a(x_a) = t_a(x_a) + x_a t_a'(x_a)$.

Las condiciones KKT del óptimo social son:
\[
\sum_{a \in p} \CMS_a(x_a^{\SO}) \geq \lambda_w \quad \forall p \in \calP_w,
\]
con igualdad si $f_p^{\SO} > 0$. Esto es un ``equilibrio de Wardrop'' con función de costo $\tilde{t}_a(x_a) = \CMS_a(x_a) = t_a(x_a) + x_a t_a'(x_a)$, que difiere de $t_a(x_a)$ siempre que $x_a > 0$ y $t_a'(x_a) > 0$.

El equilibrio de Wardrop estándar usa la función de costo $t_a$. El óptimo social corresponde a un equilibrio de Wardrop con función de costo $\CMS_a$. Dado que $\CMS_a \neq t_a$ en general, $f^{\UE} \neq f^{\SO}$ en general.
\end{proof}

\begin{example}[Red de Pigou]\label{ex:pigou}
Considérese la red más simple: dos aristas paralelas de $s$ a $d$, con demanda $d_w = 1$.
\begin{itemize}
    \item Arista 1: $t_1(x) = x$ (costo lineal, congestible).
    \item Arista 2: $t_2(x) = 1$ (costo constante, no congestible).
\end{itemize}
\textbf{Equilibrio de Wardrop:} En el equilibrio, $t_1(x_1) = t_2(x_2)$, con $x_1 + x_2 = 1$. Esto da $x_1 = 1, x_2 = 0$ y $\TC(x^{\UE}) = 1 \cdot 1 + 0 \cdot 1 = 1$.

\textbf{Óptimo social:} $\min_{x_1 \geq 0} x_1^2 + (1-x_1)$ da $x_1^* = 1/2$, con $\TC(x^{\SO}) = (1/2)^2 + (1/2) = 3/4$.

El $\PoA = 1 / (3/4) = 4/3$. Esta red es el \emph{ejemplo de Pigou} \cite{Pigou1920}, el más simple que exhibe ineficiencia del equilibrio egoísta.
\end{example}

\section{El peaje de Pigou-Knight}

\begin{theorem}[Internalización de externalidades]\label{thm:pigou_toll}
Sea $x^{\SO}$ una solución del problema del óptimo social. Definamos el \emph{peaje de Pigou} para la arista $a$ como:
\[
\tau_a^* \defeq x_a^{\SO} \cdot t_a'(x_a^{\SO}) = e_a(x_a^{\SO}).
\]
Entonces el equilibrio de Wardrop en la red con funciones de costo modificadas $\tilde{t}_a(x_a) \defeq t_a(x_a) + \tau_a^*$ coincide con $x^{\SO}$.
\end{theorem}

\begin{proof}
Con las funciones de costo modificadas $\tilde{t}_a$, el costo de la ruta $p$ en el flujo $x^{\SO}$ es:
\[
\tilde{C}_p(x^{\SO}) = \sum_{a \in p} \tilde{t}_a(x_a^{\SO}) = \sum_{a \in p} \left[ t_a(x_a^{\SO}) + x_a^{\SO} t_a'(x_a^{\SO}) \right] = \sum_{a \in p} \CMS_a(x_a^{\SO}).
\]
Las condiciones KKT del problema $(\calSO)$ evaluadas en $x^{\SO}$ establecen que:
\[
\sum_{a \in p} \CMS_a(x_a^{\SO}) \geq \lambda_w \quad \forall p \in \calP_w,
\]
con igualdad si $f_p^{\SO} > 0$. Estas son exactamente las condiciones de Wardrop para la red con costos $\tilde{t}_a$. Por tanto, $x^{\SO}$ es el equilibrio de Wardrop de la red con peajes $\tau^*$.
\end{proof}

\begin{corollary}
Para la función BPR con parámetros $\alpha$ y $\beta$, el peaje de Pigou en la arista $a$ es:
\[
\tau_a^* = \frac{\alpha \beta \, t_a^0}{c_a^\beta} (x_a^{\SO})^\beta.
\]
Para $\alpha = 0.15$, $\beta = 4$: $\tau_a^* = \dfrac{0.6 \, t_a^0}{c_a^4} (x_a^{\SO})^4$.
\end{corollary}

\begin{remark}[Significado económico]
El peaje de Pigou $\tau_a^* = e_a(x_a^{\SO})$ cobra a cada conductor exactamente la externalidad que impone al resto. Al pagar este precio, el conductor ``internaliza'' el costo externo, y su cálculo egoísta coincide con el social. El resultado es que el equilibrio descentralizado reproduce el óptimo centralizado, sin necesidad de un coordinador.

En la práctica, implementar el peaje de Pigou exacto requiere conocer $x^{\SO}$, lo cual es difícil. Los esquemas de peaje reales (como el peaje de congestión de Londres o Estocolmo) son aproximaciones a este ideal teórico.
\end{remark}
